% ---------------------------------------------------------------------------
% DRAFT SEGMENT — the cue manipulation and how the donor is chosen
% Standalone; not \input anywhere yet. Stitch into Methods later.
% Code: hrtf/modify/donor_detail.py       (the composite)
%       hrtf/processing/midline.py        (applied to the measured arc)
%       hrtf/analysis/donor_selection.py  (vanopstal_dtfs / isim / composite_isim)
%
% !! THE CODE'S ACTIVE RULE DOES NOT YET MATCH THIS TEXT. select_donor() still
%    runs the older three-measure rule (ridge slope gate + dissimilarity target
%    band + detail-strength tie-break, on un-normalised spectra). The I_sim
%    functions this section describes are implemented and validated but are not
%    yet what select_donor calls. Either switch the rule or rewrite this back.
% ---------------------------------------------------------------------------

\subsection{Spectral cue manipulation}

Participants were trained and tested with a modified version of their own
head-related transfer functions in which the coarse spectral envelope was their
own and the fine spectral detail was another listener's. For each direction and
ear, the log-magnitude spectrum was separated by cepstral truncation into an
envelope, retaining the first four coefficients, and a residual containing
everything finer. The modified spectrum was the participant's own envelope plus
a donor's residual. Own phase was retained and each direction was rescaled to
its original broadband energy, so interaural time differences and broadband
interaural level differences were unchanged and only spectral shape finer than
the envelope scale was altered.

Retaining four coefficients follows Kulkarni and Colburn
% \citet{kulkarni_role_1998}
, who smoothed transfer functions by cepstral truncation and found listeners
could no longer discriminate smoothed from original once roughly sixteen
coefficients were retained; four is well inside the range where the change is
audible and leaves the envelope too coarse to carry an elevation-dependent
pattern. Splitting the spectrum this way rather than substituting the donor's
transfer function whole was verified to decorrelate the spectral-to-spatial map
as thoroughly as a whole substitution, whereas the complementary
composite---donor envelope with own detail---left the native map intact. Detail
therefore carries the elevation cue and the envelope carries timbre, allowing
the cue to be replaced while level, head shadow and individual coloration are
retained.

The manipulation was applied to the measured median-plane transfer functions
and the full directional set was regenerated from them through the
spherical-head model used to build the unmodified set (Sec.~\ref{sec:hrir}), so
that native and modified sets share their interaural time differences exactly
and the envelope is fitted to the pinna response rather than to the pinna
response multiplied by the head shadow.

\subsection{Donor selection}

Donor transfer functions were taken from a pool of previously recorded
participants. The donor was the only aspect of the manipulation that varied
between participants, and was chosen before the participant heard any modified
stimulus, by maximizing the correspondence between the intended and the
obtained size of the spectral perturbation.

Perturbation size was quantified with the similarity index of Van Wanrooij and
Van Opstal
% \citet{vanwanrooij_relearning_2005}
, computed on directional transfer functions defined as in their Eq.~1: the
amplitude spectrum at each elevation divided by the grand average across the
elevations of the measured arc. Correlating every modified transfer function
with every unmodified one gives a matrix $C(\epsilon_1,\epsilon_2)$; for each
modified elevation the standard deviation of its correlations with all
unmodified ones was taken, and $I_\mathrm{sim}$ is the mean of those standard
deviations over elevations, computed between 4 and 18\,kHz and averaged over
ears. A high value means the modified set still discriminates elevation as the
listener's own ears do; a low value means no elevation is well matched by any.
The donor whose composite gave an $I_\mathrm{sim}$ closest to the target was
selected.

This index was used in preference to the spectral variability index
% \citep{trapeau_fast_2016}
because of what each requires of the measurement geometry. The latter is
defined on transfer functions normalized by an average over directions
distributed around the listener, which the present recordings do not sample;
the former is normalized by an average across the elevations of the median
plane, which is exactly what was measured. It is also the measure obtained in
the study whose paradigm is closest to the present one, and in that study it
predicted acute pre-adaptation elevation gain with $r^2 = 0.94$, so it is a
proxy for the acute degradation the manipulation is intended to produce.

Absolute values of $I_\mathrm{sim}$ depend on the elevation sampling, since the
index is a standard deviation across correlations, so the present values are
not directly comparable with published ones and were placed on that scale using
a common reference. Van Wanrooij and Van Opstal report $I_\mathrm{sim} = 0.5$
between a listener's own left and right ears and $0.3$ for their earmolds; the
same own-ear comparison gives $0.37$ here (31 recordings, range $0.18$--$0.51$),
and a comparison between different listeners' corresponding ears gives $0.25$
(465 pairs). The two scales therefore differ by a constant factor, and the
perturbation produced by substituting a donor's spectral detail is of the same
relative magnitude as the earmolds to which the majority of their listeners
adapted.

\subsubsection*{Donor substitution}

Selection yields a ranked list rather than a single name, so a participant for
whom the selected donor abolished rather than degraded performance could be
moved to the next donor on that list. The ranking was fixed before the
participant heard any stimulus, so the only experimenter decision was how far
down the list to move. Any substitution, its rank and the observation prompting
it were recorded, and both discarded and replacement blocks are reported.

% --- open items ------------------------------------------------------------
% CITATIONS — verified 2026-08-18 unless marked:
%   Van Wanrooij & Van Opstal (2005) J Neurosci 25(22):5413-5424,
%     doi:10.1523/JNEUROSCI.0850-05.2005. "Relearning sound localization with a
%     new ear" (SINGULAR -- the MONAURAL perturbation study). Do NOT confuse
%     with Hofman, Van Riswick & Van Opstal (1998) Nat Neurosci, "with new ears"
%     (plural), the BINAURAL mold study. I_sim definition = their Materials and
%     Methods, "DTF similarity index"; DTF definition = their Eq. 1; the 0.5 and
%     0.3 anchors = their Fig. 7A and 7C; r^2 = 0.94 = their Fig. 6 (n=6).
%     All quoted from the full text, verified.
%   Kulkarni & Colburn (1998) Nature -- cepstral smoothing, discrimination
%     threshold near 16 coefficients. [volume/pages NOT VERIFIED]
%   Trapeau et al. (2016) JASA -- VSI. Cited here only as the measure NOT used.
%     [volume/pages NOT VERIFIED]
%
% - THE TARGET VALUE IS NOT SET. The text says "closest to the target" without
%   giving one, because it should be calibrated on our own acute elevation-gain
%   data (the degradation ladder) rather than transferred from their figure --
%   their r^2 = 0.94 regression is between-subject over 6 listeners with molds
%   and on their elevation sampling, not ours. Fill in once the ladder has been
%   run on enough participants. Until then this section is incomplete.
% - Dropped from this draft at PF's direction: the donor pool is now described
%   only as "a pool of previously recorded participants". The detail-strength
%   floor, the pool membership list and the conformance/zero-padding of older
%   recordings are all out. If a reviewer asks how donors were screened, the
%   honest answer is that they were not screened on cue strength under this
%   rule -- I_sim targeting already excludes donors that would remove the cue,
%   because removing the cue drives I_sim to the bottom of the range.
% - I_sim is blind to WHERE the correlation peak sits, so it does not separate a
%   map that moved from one that vanished (Van Opstal note this: the maximum is
%   not necessarily on the diagonal). The ridge slope is retained in the code as
%   a reported check. Decide whether it appears in the paper or only as a
%   supplementary QC figure.
% - NO REPEATABILITY ESTIMATE exists for any of these measures: every arc is
%   measured once. Re-recording one subject would fix this cheaply, and two
%   recordings HAVE been remade (AS 2026-08-18, GS 2026-07-28) -- but as
%   replacements, not repeats, so they do not serve.
% - Sec.~\ref{sec:hrir} is a placeholder label -- point it at whichever section
%   describes the HRIR recording and the spherical-head azimuth expansion.
% - Regenerate if anything changes: the own-ear and between-listener I_sim
%   values and their n, and the target once set.
